\scanpage{297}
\mspage{213}

Signalons tout de suite les corollaires suivants.

\uline{Corollaire} 4.8.14. Soit $M$ une $U$-catégorie. Conditions équivalentes :

\margin{\unsure{Z|2} \uline{Pas} \uuline{prouvé} Si \unsure{finalemt}, cf 4.15.10 et 4.15.11. p.~462/463}% slanted

a) $M$ est accessible, et est équivalente à une petite catégorie.

$a')$ $M$ est accessible, et il existe un cardinal infini $\pi_0$ tel que tous les objets de $M$ soient $\pi_0$-accessible\unsure{s} \deleted{[De plus, si $\pi_0$ est \unsure{fi}} \inserted{(i.e. $M=M_{\pi_0}$)}

\deleted{$a'')$ \unsure{(b}} $b)$ $M$ est équivalente à une petite catégorie, et est Karoubienne i.e. stable par facteurs directs.

\margin{NB Ceci implique que $M$ et $M^\circ$ sont accessibles, et je conjecture que la réciproque est vraie également.}

L'équivalence de $a)$ et $a')$ est \unsure{clair|claire}, puisque $M$ admet une filtration cardinale \inserted{gén. canonique} $(M_{\pi})_{\pi\geq\pi_0}$, formée de \deleted{petites} catégories équivalentes à des petites catégories \inserted{(condition \unsure{a)} des \unsure{fcg})}, et telle que toute telle sous-catégorie soit contenue dans une des $M_{\pi}$ (\inserted{résulte de la} condition b) des fcg). Comme $M_{\pi_0}$ est équivalente à une petite catégorie et est karoubienne (étant stable par $\varinjlim$ filtrantes grandes devant $\pi_0$), il est clair aussi que $a')\Longrightarrow b)$. Reste à prouver $b)\Longrightarrow a')$, or \deleted{si $\pi_0$} cela résulte aussitôt de 4.8.13 \inserted{appliqué à $M$}, qui montre de plus qu'on peut prendre pour $\pi_0$ n'importe quel cardinal infini tel que $\pi_0\geq\operatorname{card}\operatorname{red}\deleted{\unsure{C}} M$.
% archive sheet number: 296
